\section{Líneas de trabajo futuro}
\label{sec:lineas_futuras}

El desarrollo de Voctomix~2.0 ha puesto de manifiesto varias áreas de mejora e
investigación que quedan fuera del alcance de este Trabajo de Fin de Grado pero que
constituyen líneas de trabajo futuro de interés.

\begin{itemize}

  \item \textbf{Validación con cámaras físicas en producción real.}
  La totalidad de las pruebas descritas en este trabajo se han realizado con fuentes
  sintéticas generadas mediante FFmpeg. La validación definitiva del sistema consiste
  en conectar cámaras físicas mediante tarjetas de captura o NDI y realizar una
  retransmisión en directo completa. Este escenario está previsto en el contexto del
  grupo de investigación CyberNemo de la UPM, donde Voctomix~2.0 se utilizará para
  retransmitir seminarios y presentaciones. Una producción real permitirá detectar
  posibles problemas de sincronización, \textit{jitter} de red o estabilidad que no
  se manifiestan con fuentes sintéticas de frecuencia constante.

  \item \textbf{Soporte NDI nativo como fuente de entrada.}
  Incorporar el protocolo NDI como tipo de fuente alternativo a TCP/MKV facilitaría
  la integración con cámaras de red modernas, ordenadores de presentación y
  herramientas de producción ampliamente adoptadas. La implementación requeriría
  añadir un plugin GStreamer para NDI en el contenedor de voctocore y definir un
  nuevo tipo de fuente en el fichero de configuración, manteniendo la compatibilidad
  con las fuentes TCP existentes.

  \item \textbf{Interfaz web de control.}
  Sustituir o complementar \texttt{voctogui} con una aplicación web que consuma la
  API REST del servicio de telemetría eliminaría la dependencia de GTK~3 y permitiría
  controlar el sistema desde cualquier dispositivo con navegador, incluidas tabletas
  y teléfonos móviles. Esta mejora resulta especialmente relevante en el modelo REMI,
  donde el realizador puede encontrarse geográficamente separado del servidor de
  producción.

  \item \textbf{Escalado automático en Kubernetes.}
  Kubernetes puede arrancar automáticamente nuevas instancias de un servicio cuando
  la demanda sube, mediante el mecanismo \textit{Horizontal Pod Autoscaler} (HPA).
  Una aplicación directa sería convertir la transcodificación de salida en un
  microservicio independiente: un contenedor FFmpeg que reciba la mezcla de voctocore
  y la reenvíe hacia plataformas de distribución (RTMP, SRT). Si el evento crece y
  hay que emitir a más destinos simultáneamente, Kubernetes podría lanzar nuevas
  copias de ese transcodificador de forma automática, sin interrumpir la producción
  en curso ni tocar el núcleo del mezclador.

  \item \textbf{Distribución de salida con SRT y HLS adaptativo.}
  Actualmente la señal de programa se expone como flujo TCP en bruto por el puerto
  15000. Una mejora natural sería añadir una etapa de transcodificación con FFmpeg que
  publique simultáneamente la señal en SRT para ingesta profesional de baja latencia
  y en HLS para audiencias masivas mediante CDN, integrando ambas salidas como
  servicios adicionales en el Docker Compose existente.

  \item \textbf{Integración de subtitulación automática en tiempo real.}
  La arquitectura de eventos del servicio de telemetría facilita la integración de
  un módulo de reconocimiento de voz que genere subtítulos automáticos y los publique
  como overlay de texto dinámico. Este flujo de trabajo conectaría el audio de programa
  con un servicio de transcripción (por ejemplo, Whisper de OpenAI ejecutado de forma
  local) y alimentaría el mecanismo de rotulación existente, añadiendo accesibilidad
  a las retransmisiones sin intervención manual del operador.

  \item \textbf{Panel de monitorización en tiempo real.}
  El sistema ya publica todos los eventos de producción en RabbitMQ en formato JSON.
  A partir de ahí, un consumidor intermedio podría traducir esos eventos en métricas
  numéricas y enviarlas a Prometheus, una herramienta de monitorización ampliamente
  utilizada en infraestructuras cloud~\cite{prometheus_docs}. Sobre esas métricas,
  Grafana permitiría construir un panel visual con el consumo de CPU, la latencia de
  conmutación y el estado de cada cámara en tiempo real, y configurar alertas
  automáticas que avisen al equipo si una fuente cae o el sistema supera un umbral
  de rendimiento durante la retransmisión.

\end{itemize}
